\BourbakiExerciseGroup

\begin{enumerate}
\def\labelenumi{\arabic{enumi})}
\tightlist
\item
  一个不一定交换的序幺半群 M 是一个（乘法记法写的）幺半群，配备了一个序，使得关系 \(x \leqslant y\) 蕴含 \(zx \leqslant zy\) 与 \(xz \leqslant yz\) 对所有 \(z \in M\) 。设 G 是一个不一定交换的序群。证明：
\end{enumerate}

\begin{enumerate}
\def\labelenumi{\alph{enumi})}
\item
  若 P 表示 G 中大于单位元 e 的元素之集，则 P·P = P， \(P \cap P^{-1} = \{e\}\) 且 \(aPa^{-1} = P\) 对所有 \(a \in G\) 。证明逆命题。找出使 G 成为全序群的一个条件。
\item
  如果元素 \(\sup (x,y)\) , \(\inf (x,y)\) 之一存在，则另一个也存在，并且
\end{enumerate}

\[
\sup (x, y) = x (\inf (x, y)) ^ {- 1} y = y (\inf (x, y)) ^ {- 1} x.
\]

\begin{enumerate}
\def\labelenumi{\alph{enumi})}
\setcounter{enumi}{2}
\tightlist
\item
  元素 \(\sup(x, e)\) 与 \(\sup(x^{-1}, e)\) 交换。两个互素的元素交换。d) 由 G 的不可约元生成的子群 \(G'\) 是交换的。由此推出 VI 第 18 页的定理 2 仍然成立。
\end{enumerate}

\begin{enumerate}
\def\labelenumi{\arabic{enumi})}
\setcounter{enumi}{1}
\tightlist
\item
  设 E 是一个带有合成律 \((x, y) \mapsto xy\) （不一定结合）的格，使得对所有 \(a \in E\) ，映射 \(x \mapsto ax\) 与 \(x \mapsto xa\) 都是 E 的序自同构。设 \(x_{a}\) （相应地 \(_{a}x\) ）表示由 \((x_{a})a = x\) （相应地 \(a(_{a}x) = x\) ）定义的 E 的元素，并假设对每个 \(a \in E\) ，映射 \(x \mapsto a_{x}\) 与 \(x \mapsto_{x} a\) 是从 E 的序到相反序的同构。证明：对任意的 x, y, z，
\end{enumerate}

\[
\left(z _ {\inf (x, y)}\right) \cdot x = \left(z _ {y}\right) \cdot \sup (x, y).
\]

\begin{enumerate}
\def\labelenumi{\arabic{enumi})}
\setcounter{enumi}{2}
\item
  设 \(G\) 是一个序群，其正元素之集 \(P\) 不为 0；证明 \(G\) 是无限群，且不能有最大（或最小）元素。
\item
  设 G 是一个序群，P 为 G 的正元素之集，f 为从 G 到商群 G/H 上的自然映射。要使 \(f(P)\) 使 G/H 成为一个序群，必要且充分条件是 \(0 \leqslant y \leqslant x\) 与 \(x \in H\) 蕴含 \(y \in H\) ；此时 H 称为 G 的凸子群，且 G/H 以此方式被视为一个序群。若 G 是格序的，则 G/H 为格序且 \(f(\sup(x, y)) = \sup(f(x), f(y))\) 成立的必要且充分条件是关系 \(|y| \leqslant |x|\) 与 \(x \in H\) 蕴含 \(y \in H\) ；证明这就是说 H 是一个滤过凸子群。证明一个格序群若除自身与 \(\{0\}\) 外没有其他滤过凸子群，则是全序的（考虑由 G 的两个正元素生成的滤过凸子群）；* 由此推出（一般拓扑学，V，第 25 页，习题 1）G 同构于实数的一个加法子群。*
\item
  给出一个序群的例子，其序是非离散的，且具有非零的有限阶元素（取一个适当的序群关于满足 \(P \cap H = \{0\}\) 的子群 H 的商）。
\item
  设 G 是一个序群，设 P 是其正元素之集。证明 P - P 是 G 的最大滤过子群，并且它是凸的。商上的序关系是什么？
\item
  在群 \(\mathbf{Z}\) 中，如果 \(\mathbf{P}\) 取为由 0 和整数 \(\geqslant 2\) 组成的集合，则所得的序群是滤过的但不是格序的（证明元素 \(x\) （满足 \(x \geqslant 0\) 且 \(x \geqslant 1\) ）所成的集合有两个不同的极小元素）。
\item
  设 \(x\) 是序群中的一个元素，使得 \(y = \inf(x, 0)\) 有定义；如果 \(n > 0\) 是整数，则 \(nx \geqslant 0\) 蕴含 \(x \geqslant 0\) （我们有
\end{enumerate}

\[
n y = \inf (n x, (n - 1) x, \dots , 0) \geqslant \inf ((n - 1) x, \dots , 0) = (n - 1) y);
\]

因此 \(nx = 0\) 蕴含 \(x = 0\) 。

\begin{enumerate}
\def\labelenumi{\arabic{enumi})}
\setcounter{enumi}{8}
\item
  在格序群 G 中，证明一族 \((H_{\alpha})\) 滤过凸子群之和是滤过凸子群（利用 VI，第 12 页，推论）。
\item
  证明，对每个有限序列 \((x_{i})\) ( \(1 \leqslant i \leqslant n\) )（格序群 \(G\) 的元素），
\end{enumerate}

\[
\sup \left(x _ {i}\right) = \sum_ {i} x _ {i} - \sum_ {i <   j} \inf \left(x _ {i}, x _ {j}\right) + \dots + (- 1) ^ {p + 1} \sum_ {i _ {1} <   i _ {2} <   \dots <   i _ {p}} \inf \left(x _ {i _ {1}}, \dots , x _ {i _ {p}}\right) +   + \dots + (- 1) ^ {n + 1} \inf \left(x _ {1}, \dots , x _ {n}\right).
\]

（通过基于命题 7（VI，第 10 页）的归纳论证，并利用上确界对下确界的分配性。）

\begin{enumerate}
\def\labelenumi{\arabic{enumi})}
\setcounter{enumi}{10}
\item
  设 \((x_{i})\) 是格序群 G 的 n 个元素组成的族；对每个整数 k \((0 \leqslant k \leqslant n)\) ，设 \(d_{k}\) （相应地 \(m_{k}\) ）表示 \(\binom{n}{k}\) 个具有互不相同指标的 k 项 \(x_{i}\) 之和的下确界（相应地，上确界）。证明 \(d_{k} + m_{n-k} = x_{1} + x_{2} + \cdots + x_{n}\) 。
\item
  给定格序群 \(\mathbf{G}\) ，一个子群 \(\mathbf{H}\) （\(\mathbf{G}\) 的子群）称为余格，如果对每一对 \(x, y\) （\(\mathbf{H}\) 的元素），有 \(\sup_{\mathbf{G}}(x, y) \in \mathbf{H}\) （从而等于 \(\sup_{\mathbf{H}}(x, y)\) ）。
\end{enumerate}

\begin{enumerate}
\def\labelenumi{\alph{enumi})}
\item
  如果 \(\mathbf{G} = \mathbf{Q} \times \mathbf{Q} \times \mathbf{Q}\) （其中 \(\mathbf{Q}\) 具有通常的序），则子群 \(\mathbf{H}\) （由 \((x, y, z)\) （其中 \(z = x + y\) ）组成）是一个格但不是余格。
\item
  格序群 G 的每个滤过凸（习题 4）子群 H 都是余格。
\item
  设 G 为一个格序群，H 为 G 的任意子群。定义 H' 为 H 的有限子集的下确界之集，H'' 为 H' 的有限子集的上确界之集。证明 H'' 是包含 H 的 G 的最小余格（利用 VI 第 16 页的注记）。
\end{enumerate}

\begin{enumerate}
\def\labelenumi{\arabic{enumi})}
\setcounter{enumi}{12}
\tightlist
\item
  \begin{enumerate}
  \def\labelenumii{\alph{enumii})}
  \tightlist
  \item
    一个幺半群 \(\mathbf{M}\) 称为下半格序的（或简称半格序的），如果它是一个序幺半群，如果 \(\inf(x, y)\) 对所有 \(x, y \in \mathbf{M}\) 存在，并且如果
  \end{enumerate}
\end{enumerate}

\[
\inf (x + z, y + z) = \inf (x, y) + z
\]

对所有 \(x, y, z \in \mathbf{M}\) 。证明恒等式：

\[
\inf (x, z) + \inf (y, z) = \inf (x + y, z + \inf (x, y, z))
\]

\[
\inf (x, y, z) + \inf (x + y, y + z, z + x) = \inf (x, y) + \inf (y, z) + \inf (z, x).
\]

由这些关系中的第一个，推导出 \(x \leqslant z\) 与 \(y \leqslant z\) 蕴含 \(x + y \leqslant z + \inf(x, y)\) 。

证明命题 11 及其推论（VI，第 14-15 页）在半格序幺半群中成立。b) 设 M 为半格序幺半群，N 为子幺半群，使得 \(\inf_{\mathbf{M}}(x,y) = \inf_{\mathbf{N}}(x,y)\) 对 \(x,y\in \mathbf{N}\) 成立，且 N 的元素在 M 中可逆。证明由元素 \(x - y\) （给定 \(x,y\in \mathbf{N}\) ）组成的 M 的子群 G 是格序的。

\begin{enumerate}
\def\labelenumi{\arabic{enumi})}
\setcounter{enumi}{13}
\tightlist
\item
  证明：在半格序幺半群 M 中，
\end{enumerate}

\[
\inf \left(x _ {i} + y _ {i}\right) \geqslant \inf \left(x _ {i}\right) + \inf \left(y _ {i}\right)
\]

对所有包含 n 项的有限序列 \((x_{i})\) , \((y_{i})\) 成立。推导出不等式

\[
(x + y) ^ {+} \leqslant x ^ {+} + y ^ {+}, \quad | x ^ {+} - y ^ {+} | \leqslant | x - y |
\]

在任何格序群中。

\begin{enumerate}
\def\labelenumi{\arabic{enumi})}
\setcounter{enumi}{14}
\tightlist
\item
  证明：
\end{enumerate}

\[
\left| x ^ {+} - y ^ {+} \right| + \left| x ^ {-} - y ^ {-} \right| = | x - y |
\]

在格序群中（注意 \(x - y \leqslant |x - y|\) 与 \(|x| - |y| \leqslant |x - y|\) ）。

\begin{enumerate}
\def\labelenumi{\arabic{enumi})}
\setcounter{enumi}{15}
\tightlist
\item
  证明：
\end{enumerate}

\[
n. \inf (x, y) + \inf (n x, n y) = 2 n. \inf (x, y)
\]

在半格序幺半群中（参见习题 8）。由此推出 \(\inf(nx, ny) = n \cdot \inf(x, y)\) ，如果 \(\inf(x, y)\) 是正则元。

\begin{enumerate}
\def\labelenumi{\arabic{enumi})}
\setcounter{enumi}{16}
\tightlist
\item
  设 \(x, y, z, t\) 是半格序幺半群 \(M\) 的元素，使得 \(z \geqslant 0\) 且 \(t \geqslant 0\) 。证明不等式
\end{enumerate}

\[
\inf (x + z, y + t) + \inf (x, y) \geqslant \inf (x + z, y) + \inf (x, y + t).
\]

\begin{enumerate}
\def\labelenumi{\arabic{enumi})}
\setcounter{enumi}{17}
\item
  设 \((\mathbf{G}_{\iota})_{\iota \in I}\) 是以良序集 \(I\) 为指标的全序群族；在直和 \(\mathbf{G}'\) （诸 \(\mathbf{G}_{\iota}\) 的直和）上定义一个序，取严格正元素为这样的非零 \((x_{\iota})\) ：满足 \(x_{\iota} > 0\) ，其中 \(\iota\) 为使 \(x_{\iota} \neq 0\) 的第一个指标。证明 \(\mathbf{G}'\) 是全序的。
\item
  设 \(G\) 是一个加法群，\((P_{\alpha})\) 是 \(G\) 的非空子集族，使得 \(P_{\alpha} + P_{\alpha} \subset P_{\alpha}\) 且 \(P_{\alpha} \cap (-P_{\alpha}) = \{0\}\) ；设 \(G_{\alpha}\) 表示在 \(G\) 上取 \(P_{\alpha}\) 为正元素之集所得到的序群。证明 \(P + P \subset P\) 且 \(P \cap (-P) = \{0\}\) ，其中 \(P = \bigcap_{\alpha} P_{\alpha}\) ；如果 \(H\) 是在 \(G\) 上取 \(P\) 为
\end{enumerate}

正元素之集所得到的序群，证明 H 同构于诸 \(G_{\alpha}\) 的乘积的对角子群。

\begin{enumerate}
\def\labelenumi{\arabic{enumi})}
\setcounter{enumi}{19}
\tightlist
\item
  设 \(G\) 是一个加法群，且设 \(P\) 是 \(G\) 的满足下列条件的子集：
\end{enumerate}

\begin{enumerate}
\def\labelenumi{\arabic{enumi}.}
\tightlist
\item
  \(\mathrm{P} + \mathrm{P} = \mathrm{P}\) ；2. \(\mathrm{P} \cap (-\mathrm{P}) = \{0\}\) ；3. 对每个整数 \(n > 0\) ，\(nx \in \mathrm{P}\) 蕴含 \(x \in \mathrm{P}\) （条件 (C)）。
\end{enumerate}

\begin{enumerate}
\def\labelenumi{\alph{enumi})}
\item
  如果 \(a\) 是 \(G\) 的元素使得 \(a \notin P\) ，证明存在一个子集 \(P'\) （\(G\) 的一个子集），满足条件 (C)，使得 \(P \subset P'\) 且 \(-a \in P'\) （取 \(P'\) 为由这样的 \(x \in G\) 所成的集合：存在整数 \(m > 0\) 与 \(n \geqslant 0\) 以及元素 \(y \in P\) ，使得 \(mx = -na + y\) 。
\item
  由 a) 推出，P 是 G 的满足下列条件的子集 T 的交集： \(T + T = T\) , \(T \cap (-T) = \{0\}\) , \(T \cup (-T) = G\) （即使 G 成为全序群）且 \(P \subset T\) （利用佐恩引理）。
\item
  特别地，若 G 是一个每个非零元素均无限阶的加法群，则 G 的所有满足 \(T + T = T\) , \(T \cap (-T) = \{0\}\) 且 \(T \cup (-T) = G\) 的子集 T 的交集是 \(\{0\}\) 。
\end{enumerate}

\begin{enumerate}
\def\labelenumi{\arabic{enumi})}
\setcounter{enumi}{20}
\item
  一个序群称为格可序的，如果它同构于某个格序群的子群。证明：一个序群是格可序的，当且仅当对每个整数 \(n > 0\) ，关系 \(nx \geqslant 0\) 蕴含 \(x \geqslant 0\) （为证该条件充分，利用习题 20，b) 与习题 19）。证明每个格可序群都同构于全序群乘积的一个子群。
\item
  设 G 是一个格可序群（视为 Z-模），E 是向量空间 \(G_{(Q)}\) （II，第 277 页）；证明 E 上存在唯一的与 E 的加法群结构相容的序，它诱导 G 上给定的序，并且使 E 是格可序的。
\item
  设 G 是格序群 \(Z \times Z\) （其中 Z 具有通常的序），H 是由 \((2, -3)\) 生成的 G 的凸子群；证明序群 G/H 不是格可序的（参见 VI，第 30 页，习题 4）。
\item
  设 A 是一个交换环，I 是 A 的所有理想组成的集合。则 \(\mathfrak{a}(\mathfrak{b} + \mathfrak{c}) = \mathfrak{a}\mathfrak{b} + \mathfrak{a}\mathfrak{c}\) 对 I 中所有的 a、b、c 成立（I，第 107 页）；换言之，I 连同序关系 \(a \supset b\) 与合成律 \((\mathfrak{a}, \mathfrak{b}) \mapsto \mathfrak{a}\mathfrak{b}\) 是一个半格序幺半群（VI，第 31 页，习题 13），I 中两个元素 a、b 的上确界为 \(a \cap b\) ，它们的下确界为 \(a + b\) 。
\end{enumerate}

\begin{enumerate}
\def\labelenumi{\alph{enumi})}
\tightlist
\item
  设 K 为一个域，且 \(A = K[X, Y]\) 为 K 上两个未定元的多项式环。
\end{enumerate}

\[
\mathbf {a} = (\mathbf {X}), \mathbf {b} = (\mathbf {Y}), \mathbf {c} = (\mathbf {X} + \mathbf {Y})
\]

\[
\begin{array}{c} (\mathfrak {a} \cap \mathfrak {b}) + \mathfrak {c} \neq (\mathfrak {a} + \mathfrak {c}) \cap (\mathfrak {b} + \mathfrak {c}), \\ (\mathfrak {a} + \mathfrak {b}) \cap \mathfrak {b} \neq (\mathfrak {a} \cap \mathfrak {c}) + (\mathfrak {b} \cap \mathfrak {c}), \\ (\mathfrak {a} \cap \mathfrak {b}) (\mathfrak {a} + \mathfrak {b}) \neq (\mathfrak {a} (\mathfrak {a} + \mathfrak {b})) \cap (\mathfrak {b} (\mathfrak {a} + \mathfrak {b})). \end{array}
\]

\begin{enumerate}
\def\labelenumi{\alph{enumi})}
\setcounter{enumi}{1}
\tightlist
\item
  在环 A 中，1 是 X 与 Y 的一个最大公因子，但 \((\mathbf{X}) + (\mathbf{Y}) \neq \mathbf{A}\) 。
\end{enumerate}

\begin{enumerate}
\def\labelenumi{\arabic{enumi})}
\setcounter{enumi}{24}
\tightlist
\item
  设 I 是交换环 A 的理想的半格序幺半群（习题 24）。要使有限同余式组 \(x \equiv a_i(a_i)\) 只要其中任意两个有公共解便有解（也就是说，只要
\end{enumerate}

\[
a _ {i} \equiv a _ {j} (\mathfrak {a} _ {i} + \mathfrak {a} _ {j})
\]

对每对指标 \(i, j)\) 成立），其必要且充分条件是每个合成律

\[
(a, b) \mapsto a \cap b, (a, b) \mapsto a + b,
\]

在 I 中对另一个合成律是分配的（«中国剩余定理»）。可按如下步骤进行：

\begin{enumerate}
\def\labelenumi{\alph{enumi})}
\item
  如果 \((\mathfrak{a}_1 \cap \mathfrak{a}_2) + (\mathfrak{a}_1 \cap \mathfrak{a}_3) = \mathfrak{a}_1 \cap (\mathfrak{a}_2 + \mathfrak{a}_3)\) ，并且三个同余式 \(x \equiv a_i(\mathfrak{a}_i)\) ( \(i = 1, 2, 3\) ) 中任意两个都有公共解，则这三个同余式有公共解（设 \(x_{12}\) 是 \(x \equiv a_1(\mathfrak{a}_1)\) 与 \(x \equiv a_2(\mathfrak{a}_2)\) 的一个公共解，且设 \(x_{13}\) 是 \(x \equiv a_1(\mathfrak{a}_i)\) 与 \(x \equiv a_3(\mathfrak{a}_3)\) 的一个公共解；证明同余式 \(x \equiv x_{12}(\mathfrak{a}_1 \cap \mathfrak{a}_2)\) 与 \(x \equiv x_{13}(\mathfrak{a}_1 \cap \mathfrak{a}_3)\) 有公共解）。
\item
  如果任意由三个同余式组成的组，只要其中任意两个有公共解，便有公共解，证明下列分配律成立
\end{enumerate}

\[
\mathfrak {a} + (\mathfrak {b} \cap \mathfrak {c}) = (\mathfrak {a} + \mathfrak {b}) \cap (\mathfrak {a} + \mathfrak {c}), \quad \text { and } \quad \mathfrak {a} \cap (\mathfrak {b} + \mathfrak {c}) = (\mathfrak {a} \cap \mathfrak {b}) + (\mathfrak {a} \cap \mathfrak {c}).
\]

（对于第一个，注意对所有 \(x \in (\mathfrak{a} + \mathfrak{b}) \cap (\mathfrak{a} + \mathfrak{c})\) 存在 \(y \in \mathfrak{b} \cap \mathfrak{c}\) 使得 \(y \equiv x(\mathfrak{a})\) ；对于第二个，注意对所有 \(x \in \mathfrak{a} \cap (\mathfrak{b} + \mathfrak{c})\) 存在 \(y \in \mathfrak{a} \cap \mathfrak{b}\) 使得 \(y \equiv x(\mathfrak{c})\) 。）注意由 \(a)\) 与 \(b)\) ，第二个分配律蕴含第一个（参见集合论，III，第 217 页，习题 16）。

\begin{enumerate}
\def\labelenumi{\alph{enumi})}
\setcounter{enumi}{2}
\item
  对所考虑的同余式个数作归纳，用与 a) 类似的论证证明«中国剩余定理»。
\item
  * 主理想整环情形的翻译。 *
\end{enumerate}

\begin{enumerate}
\def\labelenumi{\arabic{enumi})}
\setcounter{enumi}{25}
\item
  证明理想 (X) 在多项式环 K{[}X, Y{]} （其中 K 是一个交换域）的理想幺半群中满足 VI 第 17 页命题 14 的条件 (P)，但不是极大的。
\item
  设 A 是以 \((1, e)\) 为基的二次 Z-代数，其中 \(e^{2} = -5\) 。证明 A 是一个整环；在 A 中我们有 \(9 = 3 \cdot 3 = (2 + e)(2 - e)\) ；证明 3、\((2 + e)\) 与 \((2 - e)\) 是 A 的不可约元，但不满足 VI 第 17 页命题 14 (DIV) 的条件。
\item
  设 G 是一个格序群，其正元素之集为 P。P 的两个元素 x 与 y 称为等价的，如果与其中一个互素的任何元素也与另一个互素；此关系下的等价类称为 P 的线程；设 \(\bar{x}\) 表示包含 x 的线程。
\end{enumerate}

\begin{enumerate}
\def\labelenumi{\alph{enumi})}
\item
  设 \(\overline{a}\) 与 \(\overline{b}\) 是线程；设 \(x, x_1\) 是 \(\overline{a}\) 的元素，\(y, y_1\) 是 \(\overline{b}\) 的元素；证明如果与 \(x\) 互素的每个元素也与 \(y\) 互素，则与 \(x_1\) 互素的每个元素也与 \(y_1\) 互素。按此方式在 \(\overline{a}\) 与 \(\overline{b}\) 之间定义一个关系，记作 \(\overline{a} \geqslant \overline{b}\) ；证明这是线程之集 F 的一个序。
\item
  证明 \(\mathbf{F}\) 在此序下是一个格，并且
\end{enumerate}

\[
\inf (\bar {a}, \bar {b}) = \overline {{\inf (a , b)}}, \quad \sup (\bar {a}, \bar {b}) = \overline {{\sup (a , b)}} = \overline {{a + b}}
\]

（利用 VI 第 15 页的推论 3）。证明 F 有一个最小元素，即线程 \(\overline{0} = \{0\}\) 。

\begin{enumerate}
\def\labelenumi{\alph{enumi})}
\setcounter{enumi}{2}
\item
  称两个线程 \(\overline{a}\) , \(\overline{b}\) 互素，如果 \(\inf (\overline{a},\overline{b}) = \overline{0}\) 。证明如果 \(\overline{a}\) 与 \(\overline{b}\) 是两个 \(\neq \overline{0}\) 的线程，且若 \(\overline{a} < \overline{b}\) , 则存在一个线程 \(\overline{c}\) ，与 \(\overline{a}\) 互素，使得 \(\overline{c} < \overline{b}\) 。
\item
  称一个线程 \(\bar{m} \neq \bar{0}\) 不可约，如果它是 \(\neq \bar{0}\) 的线程所成之集的极小元素。证明一个不可约线程是全序的（如果 \(a\) 与 \(b\) 是 \(\bar{m}\) 的两个元素，考虑包含 \(a - \inf(a, b)\) 与 \(b - \inf(a, b)\) 的线程，并利用 \(b\) ）。再证明 \(\bar{m}, -\bar{m}\) 与 \(\{0\}\) 的并集是一个全序子群 \(\mathrm{H}(\bar{m})\) ，即 \(\mathbf{G}\) 的一个子群。
\item
  给定 G 的互不相同的不可约线程的族 \((\bar{m}_{\iota})\) ，证明由这些 \(\bar{m}_{\iota}\) 的并集生成的子群 H 同构于这些群 \(\mathrm{H}(\bar{m}_{\iota})\) 的直和（化归到有限族的情形，并用归纳法论证）。
\item
  证明：在非零全序群的乘积（相应地，直和）中，不可约线程的元素恰是那些除一个坐标（为正的）外其余坐标全为零的元素。由此推出，一个序群至多以一种方式表示为非零全序群的乘积（相应地，直和）（换言之，因子是唯一确定的）。
\end{enumerate}

\begin{enumerate}
\def\labelenumi{\arabic{enumi})}
\setcounter{enumi}{28}
\tightlist
\item
  序群 G 称为完全格序的，如果它是滤过的，且 G 的每个有上界的非空子集在 G 中都有上确界。下半格序幺半群称为完全半格序的，如果 \(\inf(x_{\alpha})\) 对每个非空
\end{enumerate}

的元素族 \((x_{\alpha})\) （属于 M 且有下界的）存在，并且如果

\[
\inf _ {\alpha , \beta} (x _ {\alpha} + y _ {\beta}) = \inf _ {\alpha} (x _ {\alpha}) + \inf _ {\beta} (y _ {\beta})
\]

对任意两个有下界的非空族 \((x_{\alpha})\) , \((y_{\beta})\) 都成立。完全格序群是完全半格序幺半群。

\begin{enumerate}
\def\labelenumi{\alph{enumi})}
\item
  证明：滤过群 G 是完全格序的必要且充分条件是，G 的正元素之集 P 的每个非空子集在 P 中都有下确界。
\item
  完全格序群的任意乘积仍是完全格序的。
\item
  完全格序群的每个滤过凸子群都是完全格序的。
\end{enumerate}

\begin{enumerate}
\def\labelenumi{\arabic{enumi})}
\setcounter{enumi}{29}
\tightlist
\item
  设 G 是一个序群。对 G 的每个子集 A，令 \(m(A)\) （相应地，M(A)）表示 A 在 G 中的下界（相应地，上界）所成的集合；则 \(m(A) = -M(-A)\) 。a) 关系 \(A \subset B\) 蕴含 \(m(B) \subset m(A)\) 与 \(\mathbf{M}(m(A)) \subset \mathbf{M}(m(B))\) 。
\end{enumerate}

\begin{enumerate}
\def\labelenumi{\alph{enumi})}
\setcounter{enumi}{1}
\item
  我们有 \(\mathbf{A} \subset \mathbf{M}(m(\mathbf{A}))\) 与 \(\mathbf{M}(m(\mathbf{M}(\mathbf{A}))) = \mathbf{M}(\mathbf{A})\) 。
\item
  如果 \((\mathbf{A}_{\iota})\) 是 \(G\) 的任意子集族，则 \(m\left(\bigcup_{\iota} \mathbf{A}_{\iota}\right) = \bigcap_{\iota} m(\mathbf{A}_{\iota})\) 。
\item
  每个集合 M(B)，其中 B 是 G 中有上界的非空子集，称为 G 中的主集。对 G 中任意有下界的非空子集 A，集合 M(m(A)) 是包含 A 的最小主集；将其记为 \(\langle A\rangle\) ；若 \(A\subset B\) ，则 \(\langle A\rangle\subset\langle B\rangle\) ；若 A 在 G 中有下确界 a，则 \(\langle A\rangle=M(a)\) ，我们也把它写作 \(\langle a\rangle\) 。
\item
  若 A 与 B 是 G 中两个有下界的非空子集，则 \(\langle A + B \rangle = \langle \langle A \rangle + B \rangle\) 。由此推出，在 G 的主集所成的集合 \(\mathfrak{M}(G)\) 中，映射 \((A, B) \mapsto \langle A + B \rangle\) 是一个结合且交换的合成律，对此律 \(\langle 0 \rangle = P\) 是单位元，序关系 \(A \supset B\) 与该律相容，且若 G 是滤过的，则 \(\mathfrak{M}(G)\) 在此结构下是完全半格序幺半群（习题 29）。此外，映射 \(x \mapsto \langle x \rangle\) 从 G 到 \(\mathfrak{M}(G)\) 是从序群 G 到幺半群 \(\mathfrak{M}(G)\) 的一个子群上的同构。
\item
  如果 \(\mathfrak{M}(G)\) 的一个元素 A 在此幺半群的合成律下可逆，则其逆元为 \(\mathbf{M}(-\mathbf{A}) = -m(\mathbf{A})\) （注意，若 B 是 A 的逆元，则关系 \(\mathbf{B} \subset \mathbf{M}(-\mathbf{A})\) 与 \(\mathbf{A} + \mathbf{M}(-\mathbf{A}) \subset \langle 0 \rangle\) 成立，由此得 \(\langle \mathbf{A} + \mathbf{B} \rangle = \langle \mathbf{A} + \mathbf{M}(-\mathbf{A}) \rangle\) ）。
\item
  要使 \(\mathfrak{M}(G)\) 的一个元素 A 可逆，必要且充分条件是 \(x + A \subset A\) 蕴含 \(x \geqslant 0\) （写出 \(0 \in \langle A + M(-A) \rangle\) 的一个表达式）。
\end{enumerate}

\begin{enumerate}
\def\labelenumi{\arabic{enumi})}
\setcounter{enumi}{30}
\tightlist
\item
  序群 G 称为阿基米德的，如果使得 nx（n 为正整数）所成之集有下界的元素 \(x \in G\) 只有 G 的正元素。
\end{enumerate}

\begin{enumerate}
\def\labelenumi{\alph{enumi})}
\item
  序群 G 中主集所成的幺半群 \(\mathfrak{M}(G)\) 是群，当且仅当 G 是阿基米德的（利用习题 30 的 g) 与 d)）。如果此外 G 还是滤过的，则 \(\mathfrak{M}(G)\) 是完全格序群。
\item
  由此推出，一个滤过群 G 同构于某个完全格序群的子群，当且仅当 G 是阿基米德的。
\end{enumerate}

\begin{enumerate}
\def\labelenumi{\arabic{enumi})}
\setcounter{enumi}{31}
\tightlist
\item
  \begin{enumerate}
  \def\labelenumii{\alph{enumii})}
  \tightlist
  \item
    设 G 是一个完全格序群，H 是 G 的一个子群。对群 H 的每个主子集 A，令 \(x_{\mathbf{A}}\) 表示 A 在 G 中的下确界；证明映射 \(\mathbf{A} \mapsto x_{\mathbf{A}}\) 从 \(\mathfrak{M}(\mathbf{H})\) 到 G 是单射（证明 A 恰是那些 \(y \in \mathbf{H}\) （满足 \(y \geqslant x_{\mathbf{A}}\) ）所成的集合）。
  \end{enumerate}
\end{enumerate}

\begin{enumerate}
\def\labelenumi{\alph{enumi})}
\setcounter{enumi}{1}
\item
  对 H 的每个在 H 中有下界的子集 B，令 \(\langle B\rangle\) 表示由 B 在 H 中生成的主集（\(\mathfrak{M}(H)\) 的一个元素）。如果对 H 的每个在 H 中有下界的子集 B，等式 \(\inf B = \inf \langle B \rangle\) 成立（下确界在 G 中取），证明映射 \(A \mapsto x_{A}\) 是从序群 \(\mathfrak{M}(H)\) 到 G 的一个子群上的同构（参见习题 30，e)）。
\item
  如果 G 的每个元素都是 H 的某个子集的下确界，证明对 H 的每个在 H 中有下界的子集 B，有 inf B = inf \(\langle B\rangle\) （下确界在 G 中取）。（注意，当 inf B ∈ H 时这是成立的；一般情形，考虑一个元素 \(x \in G\) ，使得 \(x + inf B \in H\) ，并利用 x 是 G 的某个子集的下确界这一事实，以及习题 30，e)。
\end{enumerate}

\(d) ^{*}\) 设 \(G\) 是完全格序群 \(\mathbf{Z} \times \mathbf{Z} \times \mathbf{R}\) 。设 \(\theta > 0\) 是一个无理数；设 \(H\) 是由 \((x, y, z)\) （满足 \(\theta(z - x) + y = 0\) ）组成的子群。证明不存在序群 \(\mathfrak{M}(H)\) 到 \(G\) 的任何子群上的、其在 \(H\) 上的限制为恒等映射的同构（注意 \(\mathfrak{M}(H)\) 同构于 \(K = Z \times R\) ；考虑 \(K\) 中由这样的元素 \(u\) 组成的子群：对每个整数 \(n > 0\) 存在 \(v \in K\) 使得 \(nv = u\) ，以及 \(G\) 的类似子群）。 *

\begin{enumerate}
\def\labelenumi{\arabic{enumi})}
\setcounter{enumi}{32}
\tightlist
\item
  \begin{enumerate}
  \def\labelenumii{\alph{enumii})}
  \tightlist
  \item
    一个全序群 G 是阿基米德的，当且仅当对 G 的任意一对元素 \(x > 0\) , \(y > 0\) ，存在整数 \(n > 0\) 使得 \(y \leqslant nx\) 。
  \end{enumerate}
\end{enumerate}

* b) 任何全序的、阿基米德的完全格序群 G 都同构于 \(\{0\}\) 、Z 或 R（排除前两种情形，在 G 中固定一个 a \textgreater{} 0，并使每个 \(x \in G\) 对应于使得 \(qx \leq pa\) 的有理数 p/q 的下确界）。

\begin{enumerate}
\def\labelenumi{\alph{enumi})}
\setcounter{enumi}{2}
\item
  由此推出，每个阿基米德全序群 G 都同构于 R 的一个子群（注意 \(\mathfrak{M}(G)\) 是全序的）。*
\item
  字典序积 \(Z \times Z\) 是全序的且不是阿基米德的。
\end{enumerate}

\begin{enumerate}
\def\labelenumi{\arabic{enumi})}
\setcounter{enumi}{33}
\item
  设 \(G = Z^{N}\) 是可数无限个全序群 Z 的乘积，并令 \(H = Z^{(N)}\) 为滤过凸子群，即 G 的各因子的直和。证明格序群 G/H（VI，第 30 页，习题 4）不是阿基米德的，尽管 G 与 H 都是完全格序的。
\item
  序群 G 中的主集 A 称为有限生成的，如果存在有限集 F 使得 \(A = \langle F \rangle\) 。我们称 G 是半阿基米德的，如果每个有限生成的主集在幺半群 \(\mathfrak{M}(G)\) 中都可逆。每个格序（相应地，阿基米德）群都是半阿基米德的（习题 31，a)）。
\end{enumerate}

\begin{enumerate}
\def\labelenumi{\alph{enumi})}
\item
  证明每个滤过的半阿基米德群都是格序的（若 \(nx \geqslant 0\) ，考虑主集 \(\langle F \rangle\) ，其中 \(F = \{0, x, \ldots, (n - 1)x\}\) ，并利用习题 30 的 g) 以及它可逆这一事实）。
\item
  * 设 K 是字典序积 \(R \times R\) ，而 G 是（通常的）乘积群 \(K \times R\) 。设 \(\theta\) 是满足 \(0 < \theta < 1\) 的无理数，并设 H 是由 \(((1, 0), 0)\) 、\(((\theta, 0), \theta)\) 与 \(((0, x), 0)\) （其中 x 遍历 R）生成的子群。证明 H 作为两个全序群的乘积的子群，不是半阿基米德的（考虑由 H 的两个元素 \(((1,0),0)\) 与 \(((\theta ,0),\theta)\) 生成的主集，并证明它不可逆）。*
\item
  设 G 是一个半阿基米德滤过群；证明 \(\mathfrak{M}(G)\) 中由诸有限生成主集及其逆元生成的子幺半群是一个格序群（参见 VI，第 31 页，习题 13，b)）。
\item
  * 设 \(\theta > 0\) 是一个无理数，并设 \(G\) 是群 \(\mathbf{Z} \times \mathbf{Z}\) ，其中正元素之集 \(P\) 取为 \((x, y)\) （满足 \(x \geqslant 0\) 且 \(y \geqslant \theta x\) ）所成之集。证明 \(G\) 是阿基米德群，但在 \(\mathfrak{M}(G)\) 中，一个非 \(\langle a \rangle\) 形式的有限生成主集的逆元不是有限生成的。*
\end{enumerate}

\begin{enumerate}
\def\labelenumi{\arabic{enumi})}
\setcounter{enumi}{35}
\tightlist
\item
  设 \(\mathbf{G}\) 是一个以 \(\mathbf{P}\) 为正元素之集的滤过群。则 \(\mathbf{G}\) 同构于形如 \(\mathbf{Z}^{(I)}\) 的群，当且仅当存在映射 \(x \mapsto d(x)\) （从 \(\mathbf{P}\) 到 \(\mathbf{N}\) ），使得对所有 \(b\) ，或者 \(b \geqslant a\) ，或者存在一个元素 \(c\) ，它属于主集 \(\langle a, b \rangle\) ，使得 \(d(c) < d(a)\) 。（证明该条件在 \(\mathbf{Z}^{(I)}\) 中成立）由映射 \(d(x_{i}) = \sum_{i} x_{i}\) （当 \(x_{i} \geqslant 0\) 对所有 \(i\) 成立时）满足。
\end{enumerate}

反之，若该条件满足，证明每个主集 \(\mathbf{A} \subset \mathbf{P}\) 具有形式 \(\langle a \rangle\) ，办法是取一个元素 \(a\) 属于 \(\mathbf{A}\) ，使得 \(d(a) \leqslant d(x)\) 对所有 \(x \in \mathbf{A}\) 成立。应用 VI 第 18 页定理 2。）
% semantic-fragments: legacy-input-seam
\relax{}
